\chapter{Morphisms and Base Change}

\section{Morphisms of Schemes}

\begin{definition}
A \emph{morphism of schemes} is a morphism of locally ringed spaces.  If \(S\)
is a scheme, an \emph{\(S\)-scheme} is a scheme \(X\) together with a morphism
\(X\to S\).  A morphism of \(S\)-schemes \(X\to Y\) is a morphism of schemes
compatible with the maps to \(S\).
\end{definition}

\begin{definition}
A morphism \(f:X\to Y\) is \emph{locally of finite type} if for every affine
open \(V=\Spec A\subset Y\) and every affine open \(U=\Spec B\subset f^{-1}V\),
the \(A\)-algebra \(B\) is finitely generated.  It is \emph{of finite type} if
it is locally of finite type and quasi-compact.
\end{definition}

\begin{proposition}
It is enough to test local finite type on affine open covers of \(X\) and
\(Y\).  More precisely, if \(Y=\bigcup_i V_i\) is an affine open cover and each
\(f^{-1}V_i\) is covered by affine opens \(U_{ij}=\Spec B_{ij}\) such that
\(B_{ij}\) is a finitely generated \(\OO_Y(V_i)\)-algebra, then \(f\) is locally
of finite type.
\end{proposition}

\begin{proof}
Let \(V=\Spec A\subset Y\) and \(U=\Spec B\subset f^{-1}V\) be affine opens.  It
suffices to prove that every point of \(U\) has a principal open neighbourhood
\(D(b)\subset U\) such that \(B_b\) is a finitely generated \(A\)-algebra.  Fix
\(x\in U\).  Choose \(i,j\) with \(x\in V\cap V_i\) and
\(x\in U\cap U_{ij}\).  After shrinking inside affine schemes, there are
principal opens \(D(a)\subset V\cap V_i\) and
\(D(b)\subset U\cap U_{ij}\) with \(f(D(b))\subset D(a)\).  The ring
\(\OO_X(D(b))=B_b\) is a localization of \(B_{ij}\), and
\(\OO_Y(D(a))=A_a\) is a localization of \(A\).  Since localizations of
finitely generated algebras are finitely generated, \(B_b\) is finitely
generated over \(A_a\), hence also over \(A\) after adjoining the image of
\(a^{-1}\), which already lies in \(B_b\).  These neighbourhoods cover \(U\).
\end{proof}

\section{Fiber Products}

\begin{definition}
Given morphisms \(X\to S\) and \(Y\to S\), a \emph{fiber product}
\(X\times_S Y\) is an \(S\)-scheme equipped with projections
\[
\begin{tikzcd}
X\times_S Y \arrow[r,"\pr_2"] \arrow[d,"\pr_1"'] & Y \arrow[d]\\
X \arrow[r] & S
\end{tikzcd}
\]
such that for every \(S\)-scheme \(T\), composition with the projections gives
a bijection
\[
    \Hom_S(T,X\times_S Y)
    \simeq
    \Hom_S(T,X)\times \Hom_S(T,Y).
\]
\end{definition}

\begin{proposition}
If \(X=\Spec B\), \(Y=\Spec C\), and \(S=\Spec A\), with structure maps
corresponding to \(A\to B\) and \(A\to C\), then
\[
    X\times_S Y \simeq \Spec(B\tensor_A C).
\]
\end{proposition}

\begin{proof}
For every affine scheme \(T=\Spec R\), morphisms \(T\to \Spec(B\tensor_A C)\)
correspond to ring maps \(B\tensor_A C\to R\).  By the universal property of the
tensor product, these are the same as pairs of ring maps \(B\to R\) and
\(C\to R\) whose restrictions to \(A\) agree.  By the affine correspondence,
this is exactly the set of pairs of morphisms \(T\to X\) and \(T\to Y\) over
S.  For a general scheme \(T\), the same bijection is obtained by checking on
an affine open cover of \(T\) and gluing, because morphisms to affine schemes
are determined by compatible maps on affine opens.  Hence
\(\Spec(B\tensor_A C)\) satisfies the required universal property.
\end{proof}

\begin{theorem}
Fiber products exist in the category of schemes.
\end{theorem}

\begin{proof}
Cover \(S\) by affine opens \(S_i\).  Cover the inverse images of \(S_i\) in
\(X\) and \(Y\) by affine opens \(X_{ij}\) and \(Y_{ik}\).  The affine products
\[
    X_{ij}\times_{S_i}Y_{ik}
\]
exist by the previous proposition.  On overlaps the affine constructions agree
because tensor product satisfies the same universal property after
localization.  Therefore the affine products glue along open subschemes.  The
glued scheme carries projections to \(X\) and \(Y\), and the universal property
is local on the source: a morphism from \(T\) to the glued scheme is equivalent
to compatible morphisms from an open cover of \(T\) into the affine product
pieces.  This gives the desired fiber product.
\end{proof}

\begin{definition}
Let \(f:X\to S\) be a morphism and let \(S'\to S\) be another morphism.  The
\emph{base change} of \(X\) to \(S'\) is
\[
    X_{S'}=X\times_S S'.
\]
If \(f:X\to Y\) is a morphism of \(S\)-schemes, its base change is the induced
morphism
\[
    f_{S'}:X\times_S S'\longrightarrow Y\times_S S'.
\]
\end{definition}

\begin{proposition}
Let \(f:X\to S\) be a morphism and let \(s\in S\).  The fiber of \(f\) over
\(s\), as a scheme, is
\[
    X_s=X\times_S \Spec \kappa(s),
\]
where \(\kappa(s)\) is the residue field of \(s\).  If \(X=\Spec B\) and
\(S=\Spec A\), and \(s\) corresponds to \(\mathfrak p\subset A\), then
\[
    X_s\simeq \Spec(B\tensor_A \kappa(\mathfrak p)).
\]
\end{proposition}

\begin{proof}
The first statement is the definition of the scheme-theoretic fiber.  The
second follows from the affine construction of fiber products and the
identification \(\Spec\kappa(\mathfrak p)\to\Spec A\) induced by
\(A\to A_{\mathfrak p}\to\kappa(\mathfrak p)\).
\end{proof}

\section{Separated Morphisms}

\begin{definition}
For a morphism \(f:X\to S\), the \emph{diagonal morphism} is the unique
\(S\)-morphism
\[
    \Delta_{X/S}:X\longrightarrow X\times_S X
\]
whose two projections to \(X\) are both the identity.  The morphism \(f\) is
\emph{separated} if \(\Delta_{X/S}\) is a closed immersion.  A scheme \(X\) is
\emph{separated} if \(X\to\Spec\ZZ\) is separated.
\end{definition}

\begin{definition}
A morphism \(i:Z\to X\) is a \emph{closed immersion} if \(i\) identifies \(Z\)
homeomorphically with a closed subset of \(X\), and the morphism
\(\OO_X\to i_*\OO_Z\) is surjective.
\end{definition}

\begin{proposition}
Every affine morphism is separated.  In particular every affine scheme is
separated.
\end{proposition}

\begin{proof}
It suffices to consider \(X=\Spec B\) over \(S=\Spec A\).  Then
\[
    X\times_S X=\Spec(B\tensor_A B).
\]
The diagonal corresponds to the multiplication map
\[
    B\tensor_A B\longrightarrow B,\qquad b\otimes b'\longmapsto bb'.
\]
This ring map is surjective because \(b\) is the image of \(b\otimes 1\).
Therefore the corresponding morphism of affine schemes is a closed immersion.
The general affine case is obtained by checking over an affine open cover of
the base.
\end{proof}

\begin{proposition}
Separatedness is stable under base change and composition.
\end{proposition}

\begin{proof}
For base change, let \(X\to S\) be separated and \(S'\to S\) any morphism.
There is a cartesian square
\[
\begin{tikzcd}
X_{S'} \arrow[r,"\Delta_{X_{S'}/S'}"] \arrow[d] &
X_{S'}\times_{S'}X_{S'} \arrow[d]\\
X \arrow[r,"\Delta_{X/S}"] & X\times_S X .
\end{tikzcd}
\]
Closed immersions are stable under base change: in the affine case this says
that a surjective ring map remains surjective after tensoring, and the general
case is local.  Hence the top diagonal is a closed immersion.

For composition, let \(X\to Y\to S\) be separated.  The diagonal
\(\Delta_{X/S}\) factors as
\[
    X \xto{\Delta_{X/Y}} X\times_Y X \longrightarrow X\times_S X.
\]
The second arrow is the base change of \(\Delta_{Y/S}\) along
\(X\times_S X\to Y\times_S Y\), hence is a closed immersion.  The first arrow
is a closed immersion by separatedness of \(X\to Y\).  A composition of closed
immersions is a closed immersion.
\end{proof}

\begin{proposition}
Let \(X\) be a scheme over \(S\).  If \(X\) is separated over \(S\), and
\(U,V\subset X\) are affine opens contained in the inverse image of an affine
open of \(S\), then \(U\cap V\) is affine.
\end{proposition}

\begin{proof}
The intersection \(U\cap V\) is the inverse image of \(U\times_S V\) under the
diagonal \(X\to X\times_S X\).  Since \(U\) and \(V\) lie over an affine open of
\(S\), the product \(U\times_S V\) is affine.  The restriction
\[
    U\cap V\longrightarrow U\times_S V
\]
is a closed immersion because it is obtained from the diagonal by base change.
A closed subscheme of an affine scheme is affine, so \(U\cap V\) is affine.
\end{proof}

\section{Proper Morphisms}

\begin{definition}
A morphism \(f:X\to S\) is \emph{universally closed} if for every morphism
\(S'\to S\), the projection \(X\times_S S'\to S'\) is a closed map on
topological spaces.  A morphism is \emph{proper} if it is separated, of finite
type, and universally closed.
\end{definition}

\begin{proposition}
Properness is stable under base change and composition.
\end{proposition}

\begin{proof}
Separatedness is stable under base change and composition by the previous
section.  Finite type is stable under base change and composition because the
corresponding statements for finitely generated algebras hold after
localization and tensor product.  Universal closedness is stable under base
change by definition: a further base change of a base change is again a base
change.  It is stable under composition because a composition of closed maps is
closed, and the same remains true after arbitrary base change.  Combining these
three facts gives the assertion.
\end{proof}

\begin{theorem}
For every ring \(A\), the structure morphism \(\PP^n_A\to\Spec A\) is proper.
\end{theorem}

\begin{proof}
The morphism is of finite type because the standard affine opens
\[
    D_+(x_i)\simeq \Spec A[x_0/x_i,\ldots,x_n/x_i]
\]
are affine spaces of finite type over \(A\).  It is separated because the
intersection of two standard affine opens is a principal open in each, hence
affine; the diagonal can therefore be checked on the standard affine cover, and
on each affine piece it is a closed immersion as in the proof that affine
morphisms are separated.

It remains to prove universal closedness.  Since projective space commutes with
base change, it suffices to prove that the projection \(\PP^n_A\to\Spec A\) is
closed.  A closed subset of \(\PP^n_A\) is locally defined by homogeneous
ideals; after replacing it by a closed subscheme, it has the form
\(\Proj S/I\) for \(S=A[x_0,\ldots,x_n]\) and a homogeneous ideal \(I\).  The
image consists of primes \(\mathfrak p\subset A\) for which
\(\Proj((S/I)\tensor_A\kappa(\mathfrak p))\) is non-empty.  This condition is
equivalent to the irrelevant ideal not being contained in the radical of
\(I(S\tensor_A\kappa(\mathfrak p))\).  By the homogeneous prime avoidance
criterion, its complement is open and is described by finitely many conditions
that some power of each \(x_i\) lies in \(I\) after localization on \(A\).  Thus
the image is closed.  The same argument after replacing \(A\) by any
\(A\)-algebra proves universal closedness.
\end{proof}

\begin{corollary}
A closed subscheme of \(\PP^n_A\) is proper over \(\Spec A\).
\end{corollary}

\begin{proof}
A closed immersion is proper: it is affine, hence separated and of finite type
when the defining ideal is finitely generated locally; it is universally closed
because it is a homeomorphism onto a closed subset after every base change.
The composition of a proper morphism with a proper morphism is proper, so a
closed subscheme of projective space is proper over \(A\).
\end{proof}

\section{Finite and Projective Morphisms}

\begin{definition}
A morphism \(f:X\to Y\) is \emph{finite} if for every affine open
\(V=\Spec A\subset Y\), the inverse image \(f^{-1}V\) is affine, say
\(\Spec B\), and \(B\) is a finite \(A\)-module.  It is \emph{affine} if the
same condition holds with no finiteness requirement on \(B\).
\end{definition}

\begin{proposition}
Finiteness is local on the target.  More precisely, if \(Y=\bigcup_i V_i\) is
an affine open cover and each \(f^{-1}V_i\to V_i\) is finite, then \(f\) is
finite.
\end{proposition}

\begin{proof}
Let \(V=\Spec A\subset Y\) be affine.  Cover \(V\) by finitely many principal
opens \(D(a_j)\) each contained in some \(V_i\); this is possible because
\(V\) is quasi-compact.  The inverse image of each \(D(a_j)\) is affine
\(\Spec B_j\), with \(B_j\) finite over \(A_{a_j}\).  The affine pieces glue
because the morphism is already affine on the cover, and their coordinate rings
are localizations of a single \(A\)-algebra \(B\).  Since a module is finite if
and only if it is finite after localization at finitely many elements
generating the unit ideal, \(B\) is finite over \(A\).  Thus \(f^{-1}V\) is
affine finite over \(V\).
\end{proof}

\begin{proposition}
Finite morphisms are closed.
\end{proposition}

\begin{proof}
It suffices to work affine.  Let \(A\to B\) be finite and let
\(Z=V(I)\subset\Spec B\).  The image of \(Z\) in \(\Spec A\) consists of primes
\(\mathfrak p\) such that there is a prime of \(B/I\) lying over
\(\mathfrak p\).  Since \(B/I\) is finite over \(A/J\), where
\[
    J=\ker(A\to B/I),
\]
the lying-over theorem for integral extensions shows that every prime of
\(A/J\) has a prime of \(B/I\) above it.  Hence the image is \(V(J)\), a closed
subset of \(\Spec A\).
\end{proof}

\begin{corollary}
Every finite morphism is proper.
\end{corollary}

\begin{proof}
A finite morphism is affine.  If \(X=\Spec B\) and \(Y=\Spec A\), then the
diagonal corresponds to the multiplication map \(B\tensor_A B\to B\), which is
surjective; hence finite morphisms are separated.  They are of finite type
because a finite module is a finitely generated algebra.  They are universally
closed because after any base change \(A\to A'\), the algebra
\(B\tensor_A A'\) is finite over \(A'\), and the preceding proposition shows
that the base-changed morphism is closed.  Thus finite morphisms are proper.
\end{proof}

\begin{definition}
A morphism \(f:X\to S\) is \emph{projective} if it factors as
\[
    X\hookrightarrow \PP^n_S\longrightarrow S
\]
for some \(n\), where the first map is a closed immersion.  It is
\emph{quasi-projective} if it factors as an open immersion into a projective
\(S\)-scheme.
\end{definition}

\begin{proposition}
Projective morphisms are proper.
\end{proposition}

\begin{proof}
The projection \(\PP^n_S\to S\) is proper by the theorem on projective space,
and closed immersions are proper.  A composition of proper morphisms is proper.
\end{proof}

\section{Flat, Smooth, and Etale Morphisms}

\begin{definition}
A morphism \(f:X\to Y\) is \emph{flat} if the local homomorphism
\[
    \OO_{Y,f(x)}\longrightarrow \OO_{X,x}
\]
is flat for every \(x\in X\).  It is \emph{faithfully flat} if it is flat and
surjective.
\end{definition}

\begin{fact}
Flatness of a module is local on source and target: an \(A\)-module \(M\) is
flat if and only if \(M_{\mathfrak p}\) is flat over \(A_{\mathfrak p}\) for
all prime ideals \(\mathfrak p\subset A\), and equivalently after localization
at a family of elements generating the unit ideal.
\end{fact}

\begin{proposition}
Flatness is stable under base change and composition.
\end{proposition}

\begin{proof}
For composition, the tensor product of flat modules is flat: if \(B\) is flat
over \(A\) and \(C\) is flat over \(B\), then \(-\tensor_A C\) is the composite
of the exact functors \(-\tensor_A B\) and \(-\tensor_B C\).  For base change,
if \(B\) is flat over \(A\), then \(B\tensor_A A'\) is flat over \(A'\) because
\[
    -\tensor_{A'}(B\tensor_A A')\simeq -\tensor_A B
\]
on \(A'\)-modules.  These affine statements localize to schemes.
\end{proof}

\begin{definition}
A morphism \(f:X\to Y\) is \emph{smooth of relative dimension \(d\)} if it is
locally of finite presentation, flat, and all geometric fibers
\[
    X_{\overline y}=X_y\tensor_{\kappa(y)}\overline{\kappa(y)}
\]
are regular schemes of dimension \(d\).  A morphism is \emph{etale} if it is
smooth of relative dimension \(0\).
\end{definition}

\begin{proposition}
Smooth and etale morphisms are stable under base change and composition.
\end{proposition}

\begin{proof}
Local finite presentation and flatness are stable under base change and
composition.  For geometric fibers, the fiber of a base change is obtained from
the original fiber by an extension of the residue field, and regularity of
geometric fibers is preserved under such extension by definition.  For
composition, the geometric fiber of \(X\to S\) over a geometric point
\(\overline s\) is a smooth family over the regular scheme
\(Y_{\overline s}\); the standard commutative algebra criterion for smooth
algebras shows it is regular of the sum of the relative dimensions.  The
relative dimension \(0\) case gives the etale statement.
\end{proof}

\begin{remark}
The differential criterion for smoothness is postponed to the chapter on
differentials.  The present definition is useful here because it behaves well
under base change without introducing differential modules prematurely.
\end{remark}

\section{Valuative Criteria}

\begin{theorem}[Valuative criterion for separatedness]
Let \(f:X\to S\) be a morphism.  Then \(f\) is separated if and only if for
every valuation ring \(R\) with fraction field \(K\), every commutative diagram
\[
\begin{tikzcd}
\Spec K \arrow[r] \arrow[d] & X \arrow[d,"f"]\\
\Spec R \arrow[r] & S
\end{tikzcd}
\]
has at most one dotted lift \(\Spec R\to X\).
\end{theorem}

\begin{proof}
If \(f\) is separated and two lifts \(g_1,g_2:\Spec R\to X\) agree on
\(\Spec K\), then the map \((g_1,g_2):\Spec R\to X\times_S X\) sends the
generic point into the diagonal.  Since the diagonal is closed, the inverse
image of the diagonal is closed in \(\Spec R\) and contains the generic point.
It is therefore all of \(\Spec R\), so \(g_1=g_2\).

Conversely, if the diagonal were not closed, there would be a point in the
closure of \(\Delta(X)\) not lying on \(\Delta(X)\).  The valuative criterion
for specialization in schemes gives a valuation ring whose generic point maps
to \(\Delta(X)\) and whose closed point maps to this boundary point.  The two
projections then give two distinct lifts agreeing generically, contradicting
the assumed uniqueness.  The specialization fact is a standard commutative
algebra application of dominating local rings by valuation rings.
\end{proof}

\begin{theorem}[Valuative criterion for properness]
Let \(f:X\to S\) be separated and of finite type, with \(S\) noetherian.  Then
\(f\) is proper if and only if for every valuation ring \(R\) with fraction
field \(K\), every commutative diagram
\[
\begin{tikzcd}
\Spec K \arrow[r] \arrow[d] & X \arrow[d,"f"]\\
\Spec R \arrow[r] & S
\end{tikzcd}
\]
has a unique dotted lift \(\Spec R\to X\).
\end{theorem}

\begin{proof}
Uniqueness is the separatedness criterion.  If \(f\) is proper, the image of
the morphism \(\Spec K\to X\) has closure finite type over \(\Spec R\), and the
closedness of \(X_R\to\Spec R\) forces the closed point of \(\Spec R\) to lie
in that closure; this gives the required lift after replacing the closure by
the scheme-theoretic image.  Conversely, suppose the lifting property holds.
To show universal closedness, it suffices by noetherian reduction to test
whether every specialization in the target of a base change lifts from the
image.  Such specializations are detected by valuation rings.  The existence
of the dotted lift says precisely that the image is stable under
specialization, hence closed in the noetherian target.
\end{proof}

\section{Fibers and Field Extensions}

\begin{definition}
Let \(X\) be a \(k\)-scheme and let \(K/k\) be a field extension.  The
\emph{extension of scalars} of \(X\) from \(k\) to \(K\) is
\[
    X_K=X\times_{\Spec k}\Spec K.
\]
If \(X=\Spec A\), then \(X_K=\Spec(A\tensor_k K)\).
\end{definition}

\begin{proposition}
Let \(X\) be a \(k\)-scheme and \(K/k\) a field extension.  Then \(K\)-points of
\(X_K\) are naturally the same as \(K\)-valued points of \(X\):
\[
    X_K(K)=X(K)=\Hom_k(\Spec K,X).
\]
\end{proposition}

\begin{proof}
By the universal property of the fiber product,
\[
    \Hom_K(\Spec K,X_K)
    \simeq
    \Hom_k(\Spec K,X)\times_{\Hom_k(\Spec K,\Spec k)}
    \Hom_k(\Spec K,\Spec K).
\]
The second factor contains the identity structure map, and the fiber product
therefore identifies with \(\Hom_k(\Spec K,X)\).
\end{proof}

\begin{proposition}
Let \(f:X\to Y\) be a morphism locally of finite type over a field \(k\).  For a
point \(x\in X\) with image \(y=f(x)\), there is a dimension inequality
\[
    \dim_x X \ge \dim_y Y+\dim_x X_y.
\]
If \(f\) is flat and \(X,Y\) are locally noetherian, then equality holds locally
in the usual dimension formula.
\end{proposition}

\begin{proof}
The inequality and the equality in the flat case are the scheme-theoretic form
of the dimension theorem for finitely generated algebras and the dimension
formula for flat local homomorphisms.  They are local on \(X\) and \(Y\), so one
reduces to a local homomorphism \(A_{\mathfrak p}\to B_{\mathfrak q}\) of
noetherian local rings essentially of finite type.  The corresponding
commutative algebra statements give the displayed formulas.
\end{proof}

\section{Exercises Used Later}

\begin{exercise}[Finite morphisms are affine]\label{ex:finite-affine}
Let \(f:X\to Y\) be finite.  Show that the inverse image of every affine open
subset of \(Y\) is affine.
\end{exercise}

\begin{solution}
This is part of the definition used in these notes.  If one defines finite
locally on an affine cover of \(Y\), then the locality proposition above shows
that the property descends to every affine open: finite algebras glue over a
finite principal cover and finiteness is local for modules.
\end{solution}

\begin{exercise}[Fibers of closed immersions]\label{ex:fiber-closed-immersion}
Let \(i:Z\hookrightarrow X\) be a closed immersion over \(S\), and let
\(S'\to S\) be any morphism.  Show that
\[
    Z\times_S S'\hookrightarrow X\times_S S'
\]
is a closed immersion.
\end{exercise}

\begin{solution}
The assertion is affine local on \(X\) and \(S\).  A closed immersion is given
by a surjection \(A\to A/I\).  After base change by \(A\to A'\), the map becomes
\[
    A'\longrightarrow (A/I)\tensor_A A'\simeq A'/IA',
\]
again a surjection.  Hence the base-changed morphism is a closed immersion.
\end{solution}

\begin{exercise}[Proper maps to affine space]\label{ex:proper-affine-functions}
Let \(X\) be proper over a field \(k\).  Show that every morphism
\(X\to\AA^1_k\) has closed image.  If \(X\) is integral and the image contains
infinitely many closed points, show that the induced rational function on \(X\)
is non-constant.
\end{exercise}

\begin{solution}
Proper morphisms are universally closed, so the image of the composition
\(X\to\AA^1_k\) is closed.  A morphism to affine space corresponds to a global
function.  If this function were constant, the image would be a single closed
point after passing to the field of constants.  Thus an image with infinitely
many closed points comes from a non-constant rational function.
\end{solution}

\begin{exercise}[Base change of projective space]\label{ex:base-change-pn}
Show that for any morphism \(S'\to S\),
\[
    \PP^n_S\times_S S'\simeq \PP^n_{S'}.
\]
\end{exercise}

\begin{solution}
Use the standard affine cover.  Over an affine open \(\Spec A\subset S\) and
\(\Spec A'\subset S'\), the inverse image of \(D_+(x_i)\) is
\[
    \Spec A[x_0/x_i,\ldots,x_n/x_i]\times_{\Spec A}\Spec A'
    \simeq
    \Spec A'[x_0/x_i,\ldots,x_n/x_i].
\]
These affine isomorphisms agree on the standard intersections, so they glue to
an isomorphism of projective spaces.
\end{solution}

\begin{exercise}[Uniqueness in the valuative criterion]\label{ex:valuative-separated}
Let \(X\to S\) be separated and let \(R\) be a DVR with fraction field \(K\).
Show directly that two \(S\)-morphisms \(\Spec R\to X\) agreeing on
\(\Spec K\) are equal.
\end{exercise}

\begin{solution}
The pair of morphisms gives \(\Spec R\to X\times_S X\).  Agreement on
\(\Spec K\) means the generic point maps into the diagonal.  Since \(X\to S\)
is separated, the diagonal is closed.  The inverse image of the diagonal is a
closed subset of \(\Spec R\) containing the generic point.  The only such
closed subset is all of \(\Spec R\), so the pair factors through the diagonal,
and the two morphisms are equal.
\end{solution}

\section{Constructible Images and Zariski Main Theorem}

\begin{definition}
A subset \(E\) of a noetherian topological space \(X\) is \emph{constructible}
if it is a finite union of locally closed subsets.
\end{definition}

\begin{theorem}[Chevalley's theorem]
Let \(f:X\to Y\) be a morphism of finite type between noetherian schemes.  If
\(E\subset X\) is constructible, then \(f(E)\subset Y\) is constructible.
\end{theorem}

\begin{proof}
The proof is noetherian induction on the target and reduces to the affine case
\(\Spec B\to\Spec A\) with \(B\) a finitely generated \(A\)-algebra.  For an
irreducible closed subset of \(\Spec B\), replace \(B\) by a domain and \(A\) by
the image domain.  Noether normalization after localizing \(A\) gives elements
of \(B\) algebraically independent over a localization of \(A\) such that \(B\)
is finite over a polynomial algebra.  The image of a non-empty open subset is
then constructible because finite maps are closed and projections from affine
space have constructible image by elimination.  Applying the induction to the
closed complements gives the result for all constructible sets.
\end{proof}

\begin{theorem}[Zariski main theorem, working form]
Let \(f:X\to Y\) be a quasi-finite separated morphism of finite type, with
\(Y\) noetherian.  Then \(f\) factors as
\[
    X\hookrightarrow \overline X\longrightarrow Y,
\]
where the first morphism is an open immersion and the second is finite.
\end{theorem}

\begin{proof}
This is stated in the form used later.  The proof starts affine locally on
\(Y\), embeds \(X\) into an affine space over \(Y\), takes the scheme-theoretic
closure in a projective compactification, and then normalizes the closure in
the finite product of residue fields coming from the quasi-finite generic
fibers.  Quasi-finiteness makes the resulting normalization finite over a
neighbourhood of each point of \(Y\), and separatedness identifies \(X\) with
an open subscheme of it by the valuative criterion.  The affine local
factorizations agree on overlaps by uniqueness of normalization, and therefore
glue.
\end{proof}

\begin{corollary}
A quasi-finite proper morphism of noetherian schemes is finite.
\end{corollary}

\begin{proof}
By Zariski main theorem, factor \(f\) as
\[
    X\hookrightarrow \overline X\to Y
\]
with the first map open and the second finite.  Since \(f\) is proper, the
image of \(X\) in \(\overline X\) is also closed: it is the image of the proper
morphism \(X\to\overline X\).  Thus \(X\) is both open and closed in
\(\overline X\).  A closed subscheme of a finite \(Y\)-scheme is finite over
\(Y\), so \(f\) is finite.
\end{proof}

\section{Blow-ups}

\begin{definition}
Let \(X\) be a scheme and let \(\II\subset\OO_X\) be a quasi-coherent ideal
sheaf.  The \emph{blow-up of \(X\) along \(\II\)} is
\[
    \operatorname{Bl}_{\II}X
    =
    \Proj\left(\bigoplus_{n\ge0}\II^n\right),
\]
where the graded algebra is the Rees algebra of \(\II\).
\end{definition}

\begin{proposition}
If \(X=\Spec A\) and \(\II=\widetilde I\), then
\[
    \operatorname{Bl}_{\II}X
    =
    \Proj\left(\bigoplus_{n\ge0}I^n\right).
\]
If \(I=(f_1,\ldots,f_r)\), the blow-up is covered by affine opens on which
\(f_i\) generates the pullback ideal.
\end{proposition}

\begin{proof}
The Rees algebra sheaf restricts on \(\Spec A\) to the graded algebra
\(\bigoplus I^n\), so the first statement is the affine definition of Proj.
The standard opens \(D_+(f_iT)\) cover the Proj because the elements \(f_iT\)
generate the positive-degree part after radical.  On \(D_+(f_iT)\), the degree
zero ring is generated over \(A\) by the ratios \(f_j/f_i\).  Thus the pullback
of \(I\) is generated by \(f_i\) on this open.
\end{proof}

\begin{theorem}[Universal property of blow-up]
Let \(\pi:\operatorname{Bl}_{\II}X\to X\) be the blow-up.  The ideal
\(\II\OO_{\operatorname{Bl}_{\II}X}\) is invertible.  If \(g:Y\to X\) is a
morphism such that \(\II\OO_Y\) is invertible, then there exists a unique
morphism \(h:Y\to\operatorname{Bl}_{\II}X\) over \(X\).
\end{theorem}

\begin{proof}
The assertion is local on \(X\), so take \(X=\Spec A\) and \(I=(f_1,\ldots,f_r)\).
On the affine chart \(D_+(f_iT)\), the pullback of \(I\) is generated by
\(f_i\), hence is invertible; these charts cover the blow-up.  Now suppose
\(I\OO_Y\) is invertible.  Locally on \(Y\), one of the generators \(f_i\) of
the pulled-back ideal generates all of \(I\OO_Y\), so the ratios \(f_j/f_i\)
are regular functions on that open.  These functions define a morphism to the
chart \(D_+(f_iT)\).  On overlaps the ratios agree, so the local morphisms glue.
Uniqueness follows because the blow-up charts are defined by precisely these
ratios.
\end{proof}

\begin{example}
The blow-up of \(\AA^2_k=\Spec k[x,y]\) at the origin is covered by two affine
charts
\[
    \Spec k[x,y/x]\qquad\text{and}\qquad \Spec k[x/y,y].
\]
The exceptional divisor is isomorphic to \(\PP^1_k\).
\end{example}
